\documentclass{article}
\input{../../lib.sty}

\title{\texttt{fn Repetitions::resolve}}
\author{Michael Shoemate}
\date{}

\begin{document}
\maketitle

Proves soundness of \texttt{fn Repetitions::resolve} in \texttt{mod.rs}.

\section{Hoare Triple}
\subsection*{Precondition}
\texttt{mean} is finite.

\subsection*{Pseudocode}
\lstinputlisting[language=Python,firstline=2,escapechar=|]{./pseudocode/repetitions_resolve.py}

\subsection*{Postcondition}
\begin{theorem}
\label{postcondition}
For any repetition specification \texttt{repetitions} and finite \texttt{mean},
\texttt{repetitions.resolve(mean)} either returns \texttt{Err(e)}, or returns a
\texttt{ResolvedRepetitions} value satisfying the proof definition in \texttt{mod.rs}.
\end{theorem}

\section{Proof}
Assume the precondition.

\begin{lemma}
\label{poisson}
If \texttt{repetitions = Poisson} and the function returns \texttt{Ok(out)}, then
\texttt{out = Poisson \{ mean \}} with nonnegative mean.
\end{lemma}

\begin{proof}
The only branch-specific check is line \ref{line:poisson-mean-guard}. If it succeeds, the function
returns exactly \texttt{Poisson \{ mean \}}.
\end{proof}

\begin{lemma}
\label{nb}
If \texttt{repetitions = NegativeBinomial \{ eta \}} and the function returns \texttt{Ok(out)},
then \texttt{eta} is finite and nonnegative, \texttt{mean > 1}, and
\texttt{out = NegativeBinomial \{ eta, x \}} where \texttt{x} is the next floating-point value
above a solution returned by \texttt{solve\_nb\_x\_from\_mean(mean, eta)}.
\end{lemma}

\begin{proof}
Lines \ref{line:nb-mean-guard}, \ref{line:eta-finite-guard}, and \ref{line:eta-sign-guard}
enforce the parameter restrictions. The only remaining fallible helper is
\texttt{solve\_nb\_x\_from\_mean} on line \ref{line:solve}; if it succeeds, the function rounds
its result upward on line \ref{line:next-up} and returns exactly the stated resolved value.
\end{proof}

\begin{proof}[Proof of Theorem~\ref{postcondition}]
By Lemma~\ref{poisson} and Lemma~\ref{nb}.
\end{proof}

\end{document}
